\documentclass[english,10pt,a4paper]{article}
\usepackage[margin=0.82in]{geometry}
\usepackage{fontspec}
\setmainfont{lmroman10-regular.otf}[BoldFont=lmroman10-bold.otf,ItalicFont=lmroman10-italic.otf,BoldItalicFont=lmroman10-bolditalic.otf]
\setsansfont{lmsans10-regular.otf}[BoldFont=lmsans10-bold.otf]
\setmonofont{lmmono10-regular.otf}[Scale=MatchLowercase]
\newlength{\LeanCharWidth}
\usepackage{amsmath,amssymb}
\usepackage[english]{babel}
\usepackage{microtype,parskip}
\usepackage{graphicx,float,caption}
\usepackage{longtable,booktabs,array,calc}
\usepackage{xcolor}
\usepackage{xurl}
\usepackage{hyperref,bookmark}
\usepackage{etoolbox,needspace}
\setcounter{secnumdepth}{-\maxdimen}
\setlength{\emergencystretch}{3em}
\providecommand{\tightlist}{\setlength{\itemsep}{0pt}\setlength{\parskip}{0pt}}
\newcounter{none}
\AtBeginEnvironment{longtable}{\Needspace{8\baselineskip}\small}
\pretocmd{\section}{\Needspace{5\baselineskip}}{}{}
\pretocmd{\subsection}{\Needspace{4\baselineskip}}{}{}
\pretocmd{\subsubsection}{\Needspace{4\baselineskip}}{}{}
\captionsetup{font=small,labelfont=bf}
\urlstyle{same}
\clubpenalty=10000\widowpenalty=10000
\DeclareMathSizes{9}{10}{7}{5}
\hypersetup{pdflang={en-US},colorlinks=true,linkcolor=blue,urlcolor=blue,citecolor=blue}
\hypersetup{pdftitle={Lean companion to Stability of the Clique Partition Gap in Split Graphs},pdfauthor={Juan Pablo Traverso Gianini}}
\title{Lean companion to Stability of the Clique Partition Gap in Split Graphs}
\author{Juan Pablo Traverso Gianini\\\small Independent researcher, Santiago, Chile\\\small\href{mailto:jtraverso@gmail.com}{jtraverso@gmail.com}\quad\href{https://orcid.org/0009-0003-6068-4096}{ORCID 0009-0003-6068-4096}}
\date{Preprint 1.0 \enspace · \enspace November 23, 2026}
\begin{document}\raggedbottom\maketitle\vspace{-1.5em}
\begin{center}\small Paper II in the series on clique partition gaps in split graphs.\end{center}
\noindent\textbf{Reserved DOI:} \url{https://doi.org/10.5281/zenodo.24452724}\\
\textbf{Official series repository:} \url{https://github.com/jtraverso/split-clique-partition-gap}\\
\textbf{Paper II repository:} \url{https://github.com/jtraverso/split-clique-partition-gap-ii}
\subsection{What the formal statement says}\label{what-the-formal-statement-says}

Preprint 1.0 uses the unrestricted integral and fractional edge-clique partition numbers from Part I. The Lean development proves the two-sided estimate (1) of Theorem 1 on the original split graph. Its input includes a residual graph, independent host neighborhoods and a proper partial assignment on actual core pairs. The parameter counting missed pairs is the cardinality of the complement of this assignment.

Lemma 4 charges the two directions separately to the unused capacities of optimal residual packings. Corollary 5 permits arbitrary crossing-edge deletions from a complete assignment, including deletions which isolate hosts. Proposition 6 bounds the change in the residual packing gap under edits on a fixed core. In Proposition 7 the compiled list bound takes the complete-core threshold. Corollary 8 bounds the unrestricted gap by twice the sum of endpoint excesses. In Theorem 9 the formal argument starts with finite convex decompositions into legal matchings and proves the unweighted product estimate and its exponential comparison. Edmonds' polytope characterization supplies those decompositions in the manuscript.

\subsection{Reproducing the checks}\label{reproducing-the-checks}

The official repository is \url{https://github.com/jtraverso/split-clique-partition-gap-ii}. Its README maps these statements to Lean declarations. The source archive distributed beside this companion contains the package under the directory \texttt{lean}. Within that directory run the following commands using elan. First-time installation requires network access.

\begin{verbatim}
lake exe cache get
lake build SplitGapII SplitGapIIAudit
lake env lean SplitGapIIAudit.lean
lake env lean AllAxioms.lean
lake env lean SplitGapII11.lean
\end{verbatim}

The package pins Lean 4.35.0-rc3 and the Mathlib revision in its Lake manifest. The first audit prints the axiom dependencies of 25 named declarations. The second scans the imported project modules and rejects any theorem using axioms outside propositional extensionality, classical choice and quotient soundness. The validation directory records the clean build and the exhaustive list.

The additional file SplitGapII11.lean checks the two list-size inequalities and the arithmetic of the two-unit jump. The two exact programs for Section 3 certify the optima in the tables. The combinatorial arguments in Propositions 5a and 6a, the general list threshold and the asymptotic applications are proved in the manuscript.

\end{document}
